\section{Reproducibility and validation}\label{app:repro}

\subsection{Archive build}

From the archive root:

\begin{lstlisting}[language=bash]
./build.sh
\end{lstlisting}

The script runs the Python unit tests, regenerates the JSON example, scans Lean files for
placeholders, compiles the report with \code{latexmk}, runs PDF preflight, and records metadata in
\code{validation/}.

Individual commands are:

\begin{lstlisting}[language=bash]
PYTHONPATH=reference python3 -m unittest -v reference/test_tapir_reference.py
PYTHONPATH=reference python3 reference/generate_demo.py
cd paper && latexmk -pdf -interaction=nonstopmode -halt-on-error tapir_ipra_framework.tex
python /home/oai/skills/pdfs/scripts/pdf_preflight.py \
  paper/tapir_ipra_framework.pdf
\end{lstlisting}

The build does not run Lean because no Lean or Lake executable was available in the artifact runtime.
To test the Lean sources, copy \code{lean/TAPIR} into the pinned ProveIt checkout, add imports in a
scratch module, and follow \code{lean/FORMALIZATION_PLAN.md}.

\subsection{PDF validation}

The generated PDF is rendered to page images at 150--200 dpi.  The shipping check requires:

\begin{itemize}
\item every page renders;
\item no clipped text, overlap, missing glyphs, or black boxes;
\item table rows remain legible;
\item code listings wrap within the text block;
\item the title, table of contents, headers, and page numbers are consistent.
\end{itemize}

The archive retains text preflight and font reports but omits the bulky render directory.  Rendered
images used for manual QA are outside the ZIP deliverable.

\subsection{Source reproducibility}

The source audit records the repository pin and the exact files inspected.  Since the branch may
advance, integration should repin and compare:

\begin{itemize}
\item \code{ConciseIPRA/Deduction.lean};
\item \code{ConciseIPRA/PR.lean};
\item \code{IPRADerivations/Core/Basic.lean};
\item \code{IPRADerivations/Core/Definitions.lean};
\item \code{IPRADerivations/Core/Search.lean};
\item \code{IPRADerivations/Core/Pairing.lean};
\item \code{IPRADerivations/Foundations/BoundedReflection.lean};
\item \code{IPRADerivations/Foundations/Termination.lean};
\item \code{ConciseIPRA/NativeClausalDiagonal.lean};
\item root README, endpoint, and unification documents.
\end{itemize}

If an API changes, update the source candidates and rerun the no-placeholder scan before making any
new verification claim.

\subsection{How to evaluate the proposal experimentally}

A useful first benchmark is not a grand theorem.  Select one existing packet with:

\begin{enumerate}
\item a small bounded checker;
\item a local certificate;
\item one explicit residual producer;
\item a final ordinary formula already expressible with current arithmetic libraries.
\end{enumerate}

Implement its local nodes as plans, order them into a schedule, and compile the root sequent.  Record:

\begin{itemize}
\item lines of new Lean;
\item number of imported generic theorems;
\item elaboration time and proof term size;
\item residual context before and after backward slicing;
\item comparison with a direct hand-written derivation.
\end{itemize}

This benchmark will determine whether the plan layer is sufficiently lightweight and where further
macros are justified.
